\documentclass{article}
\usepackage[utf8]{inputenc}
\usepackage[T1]{fontenc}
\usepackage{amsmath}
\usepackage{amsfonts}
\usepackage{amssymb}
\usepackage{enumitem}

\usepackage[margin=0.7in]{geometry} % Adjust the margins of document here.

\usepackage{tikz}                                          % Для простых рисунков в документе
\usetikzlibrary{matrix,arrows,decorations.pathmorphing,shapes.geometric,calc,snakes,backgrounds,arrows.meta}
\usepackage{xcolor}

\begin{document}
\pagestyle{plain}

\section*{Discrete Mathematics: IMC Problems 2001--2006}

\begin{center}
April 2026
\end{center}

\subsection*{Combinatorics}

\textbf{Problem 1.} (IMC 2001, Day 1, Problem 1)
Let \( n \) be a positive integer. Consider an \( n \times n \) matrix with entries \( 1, 2, \ldots, n^2 \) written in order starting from the top-left corner and moving along each row in turn left-to-right. We choose \( n \) entries of the matrix such that exactly one entry is chosen in each row and each column. What are the possible values of the sum of the selected entries?
\\

\textbf{Problem 2.} (IMC 2002, Day 2, Problem 2)
Two hundred students participated in a mathematical contest. They had 6 problems to solve. It is known that each problem was correctly solved by at least 120 participants. Prove that there must be two participants such that every problem was solved by at least one of these two students.
\\

\textbf{Problem 3.} (IMC 2003, Day 2, Problem 4)
Find all positive integers \( n \) for which there exists a family \( \mathcal{F} \) of three-element subsets of \( S = \{1, 2, \ldots, n\} \) satisfying:
\begin{enumerate}
    \item for any two different elements \( a, b \in S \), there exists exactly one \( A \in \mathcal{F} \) containing both \( a \) and \( b \);
    \item if \( a, b, c, x, y, z \) are elements of \( S \) such that \( \{a, b, x\}, \{a, c, y\}, \{b, c, z\} \in \mathcal{F} \), then \( \{x, y, z\} \in \mathcal{F} \).
\end{enumerate}
\\

\textbf{Problem 4.} (IMC 2004, Day 1, Problem 1)
Let \( S \) be an infinite set of real numbers such that \( |s_1 + s_2 + \cdots + s_k| < 1 \) for every finite subset \( \{s_1, s_2, \ldots, s_k\} \subset S \). Show that \( S \) is countable.
\\

\textbf{Problem 5.} (IMC 2005, Day 1, Problem 2)
For an integer \( n \geq 3 \) consider the sets
\begin{align*}
S_n &= \{(x_1, \ldots, x_n) : \forall\, i,\; x_i \in \{0,1,2\}\}, \\
A_n &= \{(x_1, \ldots, x_n) \in S_n : \forall\, i \leq n-2,\; |\{x_i, x_{i+1}, x_{i+2}\}| \neq 1\}, \\
B_n &= \{(x_1, \ldots, x_n) \in S_n : \forall\, i \leq n-1,\; (x_i = x_{i+1} \Rightarrow x_i \neq 0)\}.
\end{align*}
Prove that \( |A_{n+1}| = 3 \cdot |B_n| \).
\\

\subsection*{Number Theory and Algebra}

\textbf{Problem 6.} (IMC 2001, Day 1, Problem 2)
Let \( r, s, t \) be positive integers which are pairwise relatively prime. If \( a \) and \( b \) are elements of a commutative multiplicative group with unity element \( e \), and \( a^r = b^s = (ab)^t = e \), prove that \( a = b = e \). Does the same conclusion hold if \( a \) and \( b \) are elements of an arbitrary non-commutative group?
\\

\textbf{Problem 7.} (IMC 2002, Day 2, Problem 3)
For each \( n \geq 1 \) let
\[
a_n = \sum_{k=0}^{\infty} \frac{k^n}{k!}, \qquad b_n = \sum_{k=0}^{\infty} (-1)^k \frac{k^n}{k!}.
\]
Show that \( a_n \cdot b_n \) is an integer.
\\

\textbf{Problem 8.} (IMC 2003, Day 1, Problem 2)
Let \( a_1, a_2, \ldots, a_{51} \) be non-zero elements of a field. We simultaneously replace each element with the sum of the 50 remaining ones, obtaining \( b_1, \ldots, b_{51} \). If this new sequence is a permutation of the original one, what can be the characteristic of the field?
\\

\textbf{Problem 9.} (IMC 2006, Day 1, Problem 2)
Find the number of positive integers \( x \) satisfying the following two conditions:
\begin{enumerate}
    \item \( x < 10^{2006} \);
    \item \( x^2 - x \) is divisible by \( 10^{2006} \).
\end{enumerate}
\\

\textbf{Problem 10.} (IMC 2006, Day 2, Problem 5)
Prove that there exists an infinite number of relatively prime pairs \( (m, n) \) of positive integers such that the equation
\[
(x + m)^3 = nx
\]
has three distinct integer roots.

\end{document}
